% ===== CHƯƠNG 1: MỞ ĐẦU =====

\chapter{Mở Đầu}

\section{Lý do chọn đề tài}

Bài toán tìm đường ngắn nhất trên đồ thị có trọng số là một trong những vấn đề cơ bản nhất của Khoa học Máy tính và Trí tuệ Nhân tạo, với ứng dụng rộng rãi trong đời sống thực tế. Các thuật toán pathfinding được sử dụng trong hệ thống định vị GPS, lập kế hoạch tuyến đường giao thông, trò chơi điện tử, robot navigation, và nhiều lĩnh vực khác.

Thành phố Thượng Hải sở hữu một trong những hệ thống tàu điện ngầm (Metro) lớn nhất thế giới với hơn 400 ga và 20 tuyến, tạo ra một đồ thị phức tạp lý tưởng để kiểm thử các thuật toán. Việc xây dựng một ứng dụng thực tế để tìm đường tối ưu giữa hai ga metro không chỉ giúp hiểu rõ lý thuyết mà còn minh hoạ cách áp dụng các thuật toán vào bài toán thực tiễn.

Ngoài ra, đề tài này cho phép so sánh trực tiếp bốn thuật toán tìm kiếm khác nhau --- Uniform Cost Search, Greedy Best-First Search, A*, và Dijkstra --- trên cùng một bộ dữ liệu, từ đó rút ra các kết luận về hiệu năng, tính tối ưu, và sự phù hợp của mỗi thuật toán trong các tình huống khác nhau.

\section{Mục tiêu của đề tài}

Báo cáo này có các mục tiêu chính sau:

\begin{enumerate}
  \item \textbf{Nghiên cứu và phân tích} bốn thuật toán tìm kiếm: hiểu rõ nguyên lý hoạt động, độ phức tạp thời gian và không gian, tính hoàn chỉnh (completeness) và tính tối ưu (optimality) của từng thuật toán.

  \item \textbf{So sánh hiệu năng} của các thuật toán trên dữ liệu thực: đo lường số node được khám phá, chi phí đường đi, thời gian thực thi để rút ra những phát hiện về khi nào nên dùng thuật toán nào.

  \item \textbf{Xây dựng ứng dụng web thực tế} cho phép người dùng tìm đường giữa hai ga metro theo nhiều chế độ khác nhau, bao gồm chọn từ danh sách và click trực tiếp trên bản đồ.

  \item \textbf{Phát triển kỹ năng lập trình} về kiến trúc backend-frontend, xử lý yêu cầu HTTP, trực quan hoá dữ liệu trên bản đồ, và viết test case toàn diện.

  \item \textbf{Làm việc với dữ liệu thực} từ OpenStreetMap, hiểu những thách thức của xử lý dữ liệu địa lý, và xử lý các edge case như dữ liệu lỗi, điểm ngoài phạm vi, v.v.
\end{enumerate}

\section{Phạm vi đề tài}

Báo cáo này có các giới hạn và giả định sau:

\begin{itemize}
  \item \textbf{Dữ liệu:} Sử dụng dữ liệu mạng lưới metro Thượng Hải từ OpenStreetMap gồm 411 ga và 20 tuyến. Đây là dữ liệu sẵn có từ ngày cập nhật, không được cập nhật real-time.

  \item \textbf{Đồ thị:} Đồ thị là vô hướng (undirected), có trọng số dương, biểu diễn mạng metro nơi có thể đi cả hai chiều giữa các ga liền kề.

  \item \textbf{Chi phí cạnh:} Chi phí được ước lượng dựa trên công thức (khoảng cách / 35 km/h + phụ phí dừng), không dùng dữ liệu thời gian chạy thực tế do OSM không cung cấp.

  \item \textbf{Giới hạn địa lý:} Khi tìm ga gần nhất từ một điểm bất kỳ, hệ thống chỉ chấp nhận những điểm nằm trong bán kính 50 km từ ga gần nhất, nếu vượt quá sẽ báo lỗi.

  \item \textbf{Không tính:} Báo cáo không xét các yếu tố như giờ cao điểm, đóng cửa tuyến, độ trễ thực tế, hoặc thời gian đợi giữa các tuyến (transfer time). Đây là những yếu tố để mở rộng trong tương lai.

  \item \textbf{Một cặp (start, goal)} trong mỗi lần chạy. Không xét bài toán tìm kiếm toàn bộ cây đường ngắn nhất từ một điểm đến tất cả điểm khác (mặc dù Dijkstra có khả năng này).
\end{itemize}

\section{Phương pháp nghiên cứu}

\subsection{Phương pháp lý thuyết}

Sử dụng tài liệu học tập từ môn Nhập môn AI và các sách chuyên khảo (AIMA, CLRS) để nắm vững lý thuyết các thuật toán, bao gồm:
\begin{itemize}
  \item Đọc và phân tích pseudocode từ tài liệu.
  \item Hiểu rõ điều kiện admissibility của heuristic, tính hoàn chỉnh và tối ưu.
  \item Phân tích độ phức tạp thời gian và không gian.
\end{itemize}

\subsection{Phương pháp thực nghiệm}

Xây dựng ứng dụng và chạy các thí nghiệm:
\begin{itemize}
  \item Cài đặt đầy đủ 4 thuật toán bằng Python.
  \item Viết 29 test case bao gồm unit test (kiểm thử từng hàm) và integration test (kiểm thử toàn flow HTTP).
  \item Chạy benchmark trên những cặp ga khác nhau để đo explored node count, chi phí, thời gian thực thi.
  \item So sánh kết quả giữa các thuật toán.
\end{itemize}

\subsection{Phương pháp thiết kế}

Sử dụng quy trình thiết kế phần mềm chuẩn:
\begin{itemize}
  \item Phân tích yêu cầu (functional + non-functional).
  \item Thiết kế kiến trúc 3-tầng: Presentation, Business Logic, Data.
  \item Thiết kế API RESTful với JSON.
  \item Thiết kế giao diện người dùng (UI) responsive, sử dụng Leaflet cho bản đồ.
\end{itemize}

\section{Cấu trúc báo cáo}

Báo cáo gồm 5 chương chính:

\begin{description}
  \item[Chương 1 --- Mở đầu] Nêu lý do chọn đề tài, mục tiêu nghiên cứu, phạm vi, phương pháp, và cấu trúc báo cáo.

  \item[Chương 2 --- Cơ sở lý thuyết] Trình bày chi tiết bài toán tìm đường, biểu diễn đồ thị, phân tích 4 thuật toán (UCS, Greedy, A*, Dijkstra), hàm heuristic, và so sánh lý thuyết.

  \item[Chương 3 --- Phân tích và thiết kế] Liệt kê yêu cầu chức năng và phi chức năng, trình bày kiến trúc hệ thống, mô hình dữ liệu, thiết kế API, thiết kế giao diện.

  \item[Chương 4 --- Thực hiện và kết quả] Mô tả công nghệ sử dụng, cấu trúc mã nguồn, chi tiết cài đặt các thuật toán, kết quả kiểm thử, kết quả benchmark so sánh.

  \item[Chương 5 --- Kết luận] Tóm tắt những gì đạt được, nêu hạn chế hiện tại, và đề xuất hướng phát triển tương lai.
\end{description}

Ngoài ra, các phụ lục chứa hướng dẫn cài đặt và chạy ứng dụng, cũng như các tài liệu tham khảo.
